%% =====================================================================
%%  T-06-05  The Three Channels
%%  -------------------------------------------------------------------
%%  One idea per file. Core is mandatory. Slots in canonical order.
%%  Max 700 words. No layout commands. No filler. New source material
%%  for this entry goes in sources/<BOOK>/T-06-05.tex -- never by
%%  rewriting this file (docs/RULES.md R0.1).
%% =====================================================================
%% @kind: topic
%% @id: T-06-05
%% @title: The Three Channels
%% @subtitle: Words, tone and face --- and the yes that only travels on one of them
%% @chapter: c06
%% @order: 50
%% @status: stable
%% @sources: B5
%% @updated: 2026-09-19
%% @allow-rewrite: B5 ch. 8-10 ingest 2026-09-19 -- committed scaffold replaced by the written entry (planned -> stable lifecycle)

\topic{T-06-05}{The Three Channels}{Words, tone and face --- and the yes that only travels on one of them}

\begin{core}
A message travels on three channels at once: the words, the tone of voice, and the face and body. People rehearse the words and not the other two, so a managed statement leaks through tone and face. A yes that arrives in the words but not in the other channels is not agreement yet --- and pressing for it produces the counterfeit kind.
\end{core}
\begin{mechanism}
The classical figure is 7--38--55, from studies on communicating feelings and attitudes: about seven per cent of such a message rides on the words, thirty-eight on tone, fifty-five on face and body. The arithmetic is contested; the claim that survives is that the non-verbal channels carry most of the emotional content and are far harder to script. Lying research finds the same shape: tells cluster in the unrehearsed channels. When the channels disagree, the words are the managed part and the leakage is the truth. The yes taxonomy (counterfeit, confirmation, commitment) pairs with the rule of three: an agreement re-confirmed three times, framed differently, is triply hard to counterfeit.
\end{mechanism}
\begin{conditions}
Strongest in face-to-face and voice high-stakes agreement, where all three channels are live and the stakes make rehearsal expensive. Weakens sharply over text, where two of three channels are gone and the check becomes the writing itself: specificity, consistency, the two-tellings test (T-06-02). Quote the ratio as a pointer, not arithmetic: it comes from feeling-and-attitude studies.
\end{conditions}
\begin{application}
  \item Watch the face in the gap before the answer settles (T-07-03); weigh the words after.
  \item Listen for the mismatch: enthusiasm in the vocabulary, flatness in the voice, agreement in the words and retreat in the posture.
  \item Summarise what they agreed to and watch the reply's channels, not its content.
  \item For any consequential yes, run the rule of three: get the same agreement three times in the conversation, each time framed differently.
  \item Over text, move the check into the record: exact numbers and dates, expensive to fake, easy to re-check.
\end{application}
\begin{feedback}
  \item Landing: the second and third confirmations come faster and more specific --- commitment is real.
  \item Landing: a gently named mismatch (\emph{you sound less sure than you sound}) is met with relief, not annoyance.
  \item Failing: the re-ask shifts the story's specifics --- invented accounts regenerate (T-06-02).
  \item Failing: irritation at a harmless third confirmation. The yes was counterfeit and the check found it.
\end{feedback}
\begin{failure}
Read as a checklist, the channels produce false positives: crossed arms mean a cold room, a paused answer means a careful person. Misreading leakage as arithmetic rather than evidence makes the reader confident and wrong. There is also a training effect: run the checks visibly and the counterpart learns to manage tone and face as well as words, which produces the smooth, dead consistency of a rehearsed script --- the tell that used to work is now performed.
\end{failure}
\begin{countermeasures}
  \item Assume your own leakage. A yes you do not mean will show in your voice before it shows in your diary; that is the argument for not giving one.
  \item Separate comfort from agreement in yourself: liking someone is not a reason in the matter (T-22-01).
  \item When a counterpart triple-confirms your yes, notice the technique and ask which part they doubt --- their doubt is the real negotiation.
  \item Buy time before big agreements. Leakage scales with pressure, and pressure is cheapest to remove.
\end{countermeasures}

\sources{B5}
\seesources
\seealso{T-06-01, T-06-02, T-07-03}
